\section{Set size and the recruitment of validity-graded attention}
\label{sec:setsize}

A central prediction of any account in which attention allocates a limited resource is
that the \emph{benefit} of knowing where to look should grow with the number of competing
locations \citep{palmer1995psychophysics, mazyar2012does, eckstein2013visual}. With four
locations, a perfectly valid cue removes a modest amount of spatial uncertainty; with nine,
it removes much more, and an agent that exploits the cue should therefore recruit
reliability-graded attention more strongly at the larger set size. We use the nine-stimulus
variant (Section~\ref{sec:env-setsize}) to test this set-size $\times$ validity interaction
in the coupled architecture.

\paragraph{Design and prediction.}
The nine-stimulus environment cues one of nine locations and signals its reliability through
five proportion levels $\{0, 0.25, 0.50, 0.75, 1.00\}$, with an uninformative cue
(proportion $0$, uniform over all nine locations) as a within-task baseline and a uniform
reward structure that removes the value manipulation, isolating validity. The architecture is
the coupled cross-talk model of Section~\ref{sec:configs} applied unchanged. We make two
quantitative predictions. First, the \emph{behavioural} cueing benefit --- the advantage for a
change at the cued location, and its growth across proportion levels --- should be
\emph{larger} at nine stimuli than at four. Second, and more diagnostically, the
\emph{attentional} recruitment should be more strongly graded by proportion: the attention the
model allocates to the cued location at the moment of change should rise more steeply with cue
reliability when there are nine competing locations than when there are four. In short, the
prediction is that the nine-stimulus task uses proportion-graded attention \emph{more} than the
smaller-set-size tasks.

\paragraph{Status and analysis plan.}
We report this as a designed experiment and an open quantitative prediction rather than a
completed result: at the time of writing the nine-stimulus run has not been trained to
convergence with a logged, analysable model, and no per-reliability attention analysis exists
(\TODO{set-size $\times$ validity result --- to be filled when the canonical sweep completes}).
The analysis is specified in advance. Using matched runs at set sizes $K\in\{2,4,9\}$ under one
architecture and training recipe, we will (i) fit psychometric functions split by cue
reliability and cued/uncued location and extract the reliability-graded validity benefit at each
$K$; and (ii) read the cross-attention weights at the change frame, compute the attention
allocated to the cued versus uncued locations as a function of reliability, and test whether the
reliability slope of cued attention increases with $K$. A positive interaction --- steeper
reliability-graded attention and a larger cueing benefit at $K=9$ --- would provide a parametric,
within-model demonstration that the cue is recruited in proportion to the uncertainty it
resolves, and a corresponding prediction for set-size manipulations in primate cueing
experiments \citep{mayo2016graded, eckstein2013visual}.
